%% kapa: exam=1 round=23 year=2012 period=1 form=A subject=civil_law subject_ko=민법 seq=17 number=17
%% source: Q-Net 감정평가사 기출문제 (cst00309) artlSeq=5052578 file=raw/1008/5052578/2012년 제23회 감정평가사 제1차시험 기출문제(1교시 A형).pdf
%% conversion: kordoc 4.24.51 -> tools/gen1.py (auto); review=proofread-sonnet; flags=none
\begin{question}{17}
의사표시의 효력발생시기에 관한 설명으로 옳지 않은 것은? (다툼이 있으면 판례에 의함)
\begin{choices}
  \item 무권대리인의 상대방이 상당한 기간을 정하여 본인에게 추인 여부를 최고한 경우, 본인이 그 기간 내에 확답을 발하지 않으면 추인을 거절한 것으로 본다.
  \item 공시송달을 함에 있어서 표의자는 상대방을 알지 못하거나 그 소재를 알지 못하는 데 과실이 없어야 한다.
  \item 미성년자에 대한 의사표시의 도달을 법정대리인이 안 후에는 표의자가 의사표시의 도달을 주장할 수 있다.
  \item 보통우편의 방법으로 의사표시가 발송된 경우에는 그 의사표시가 상당한 기간 내에 도달한 것으로 추정된다.
  \item 승낙의 의사표시를 발송한 후 승낙자가 사망하더라도 승낙의 효력에는 영향을 미치지 않는다.
\end{choices}
\end{question}
\answer{4}
